\Article{64}
Starting from the last rule obtained in the preceding article, we find
\[
\begin{array}{c|l}
\text{Number}&\text{Character}\equiv\\\hline
-1+i&\frac18(aa+2ab-bb-1)\\
-1-i&\frac18(-aa+2ab+bb+1)\\
+1-i&\frac18(aa+2ab+3bb-1)
\end{array}
\]

This follows at once from the fact that the character of $i$ is $\frac14(aa+bb-1)$, while that of $-1$ is $\frac12(aa+bb-1)\equiv\frac12bb$, since $aa-1$ always has the form $8n$. These four rules, although so far obtained by examining examples, are evidently connected in such a way that, as soon as the proof of one has been completed, the other three will have been proved at the same time. It is scarcely necessary to point out that these rules too suppose only that $a$ is odd and $b$ even.

If we are willing to use formulas restricted to primary moduli, we can use the following form:
\[
\begin{array}{c|l}
\text{Number}&\text{Character}\equiv\\\hline
-1+i&\frac14(-a-b+1-bb)\\
-1-i&\frac14(a-b-1+bb)\\
+1-i&\frac14(-a-b+1+bb)
\end{array}
\]

The simplest formulas result if, as we did at the beginning of our examination of examples, we distinguish moduli of the first and second kinds. The characters are
\[
\begin{array}{c|c|c}
\text{Number}&\text{For moduli of the first kind}&\text{For moduli of the second kind}\\\hline
-1+i&\frac14(-a-b+1)&\frac14(-a-b-3)\\
-1-i&\frac14(a-b-1)&\frac14(a-b+3)\\
+1-i&\frac14(-a-b+1)&\frac14(-a-b+5)
\end{array}
\]

\Article{65}
For the number $-1+2i$, to which we now proceed, we shall likewise use the same distinction between those moduli $a+bi$ for which $a\equiv1,b\equiv0$ and those for which $a\equiv3,b\equiv2$. The table in Article~62 shows that, with respect to that number, the following correspondences hold:
\SourceLeaf{147}%
\[
\begin{array}{c|l}
\text{Character}&\text{Moduli of the first kind}\\\hline
0&-3+8i,\ +5-8i,\ +9+4i,\ -11+4i\\
1&+1+4i,\ +5-4i,\ -7,\ -3-8i\\
2&+1-4i,\ +5+8i,\ -7-8i,\ -11\\
3&-3,\ +5+4i,\ +9-4i,\ -7+8i,\ -11-4i
\end{array}
\]

Reducing each of these moduli to its absolutely least residue modulo $-1+2i$, we observe that all those to which character $0$ corresponds are $\equiv1$; those to which character $1$ corresponds are $\equiv i$; those whose character is $2$ become $\equiv-1$; finally, all those whose character is $3$ become $\equiv-i$. But the characters of $1,i,-1,-i$ for the modulus $-1+2i$ are themselves $0,1,2,3$, respectively. Thus, in all these $17$ examples, the character of $-1+2i$ with respect to a modulus $a+bi$ of the first kind is identical with the character of that number with respect to the modulus $-1+2i$.

Likewise, with the help of the table, we find the correspondences
\[
\begin{array}{c|l}
\text{Character}&\text{Moduli of the second kind}\\\hline
0&+3+2i,\ -5-2i,\ -1+10i,\ -1-10i,\ +11+6i\\
1&+3-2i,\ -1+6i,\ -5-6i,\ +7+10i,\ +7-10i\\
2&-5+2i,\ -1-6i,\ +7-2i\\
3&-1-2i,\ +7+2i,\ -5+6i,\ +3+10i,\ +3-10i,\ +11-6i
\end{array}
\]

Reducing these moduli to their least residues modulo $-1+2i$, all those to which the characters $0,1,2,3$ respectively correspond are found congruent to $-1,-i,+1,+i$. But when, conversely, $-1+2i$ is taken as the modulus, these very numbers have the characters $2,3,0,1$, respectively. Consequently, in all these $19$ examples, the character of $-1+2i$ with respect to a modulus of the second kind differs by two units from the character of that number with respect to $-1+2i$ taken as the modulus.

Moreover, it is readily seen that entirely similar results will hold for $-1-2i$.

\Article{66}
For the number $-3$, we omit the distinction between moduli of the first and second kinds, since the result shows that it is superfluous here. Thus the correspondences are
\SourceLeaf{148}%
\[
\begin{array}{c|l}
\text{Character}&\text{Moduli}\\\hline
0&\begin{gathered}-1+6i,\ -1-6i,\ -7,\ -5+6i,\\-5-6i,\ -11,\ 11+6i,\ 11-6i\end{gathered}\\[3pt]
1&\begin{gathered}-1-2i,\ 1-4i,\ -5+2i,\ 5+4i,\ 7+2i,\\5-8i,\ -1+10i,\ -7-8i,\ -11-4i,\ 7-10i\end{gathered}\\[3pt]
2&\begin{gathered}3+2i,\ 3-2i,\ -3+8i,\ -3-8i,\\9+4i,\ 3+10i,\ 3-10i\end{gathered}\\[3pt]
3&\begin{gathered}-1+2i,\ 1+4i,\ -5-2i,\ 5-4i,\ 7-2i,\\5+8i,\ -1-10i,\ -7+8i,\ -11+4i,\ 7+10i\end{gathered}
\end{array}
\]

\Needspace{12\baselineskip}
Reducing these moduli to their least residues modulo $3$, we see that those to which character $0$ corresponds are partly $\equiv1$, partly $\equiv-1$; those whose character is $1$ become either $\equiv1-i$ or $\equiv-1+i$; those whose character is $2$ become either $\equiv i$ or $\equiv-i$; finally, those to which character $3$ belongs are either $\equiv1+i$ or $\equiv-1-i$. From this examination of examples we therefore conclude that the character of $-3$ for a modulus that is a prime number primary among its associates is identical with the character of that very number when $3$, or equivalently $-3$, is considered as the modulus.

\Article{67}
Examining examples similarly for other prime numbers, we find that $3\pm2i,-1\pm6i,7\pm2i,-5\pm6i$, etc., give theorems similar to the one reached in Article~65 for $-1+2i$, whereas $1\pm4i,5\pm4i,-3\pm8i,5\pm8i,9\pm4i$, etc., behave in the same way as $-3$. The examination of examples thus leads to a most elegant theorem,. which, by analogy with the theory of quadratic residues in the arithmetic of real numbers, may be called the \textsc{Fundamental Theorem} of the theory of biquadratic residues, namely:

\emph{Let $a+bi,a'+b'i$ denote distinct prime numbers, primary among their respective associates, that is, congruent to unity modulo $2+2i$. The biquadratic character of $a+bi$ with respect to the modulus $a'+b'i$ will be identical with the character of $a'+b'i$ with respect to the modulus $a+bi$, if either both numbers $a+bi,a'+b'i$, or at least one of them, belong to the first kind, that is, are congruent to unity modulo $4$. On the other hand, those characters will differ from one another by two units if neither number $a+bi,a'+b'i$ belongs to the first kind, that is, if each is congruent modulo $4$ to $3+2i$.}

\SourceLeaf{149}%
Yet despite the great simplicity of this theorem, its proof must be counted among the most deeply hidden mysteries of higher arithmetic. At least as matters now stand, it can be unravelled only through the most subtle investigations, which would go far beyond the limits of the present paper. We therefore reserve the publication of this proof, as well as the development of the connection between this theorem and those that we had begun to establish by examining examples at the beginning of this paper, for a third paper. Nevertheless, by way of conclusion, we shall give here what is required to prove the theorems proposed in Articles~63.~64.

\Needspace{15\baselineskip}
\Article{68}
We begin with prime numbers $a+bi$ for which $b=0$ (the third species in Article~34). Here, for the number to be primary among its associates, $a$ must therefore be a negative real prime number of the form $-(4n+3)$, for which we shall write $-q$, such as $-3,-7,-11,-19$, etc. Denoting by $\lambda$ the character of $1+i$, with that number taken as the modulus, we must have
\[
i^\lambda\equiv(1+i)^{\frac14(qq-1)}\equiv2^{\frac18(qq-1)}\pdot i^{\frac18(qq-1)}\pmod q
\]
But it is known that $2$ is a quadratic residue or a quadratic nonresidue of $q$ according as $q$ has the form $8n+7$ or the form $8n+3$. From this we conclude that, in general,
\[
2^{\frac12(q-1)}\equiv(-1)^{\frac14(q+1)}\equiv i^{\frac12(q+1)}\pmod q
\]
and consequently, raising to the power with exponent $\frac14(q+1)$,
\[
2^{\frac18(qq-1)}\equiv i^{\frac18(q+1)^2}\pmod q
\]
The preceding equation therefore takes the form
\[
i^\lambda\equiv i^{\frac18(q+1)^2+\frac18(qq-1)}\equiv i^{\frac14(qq+q)}\pmod q
\]
whence it follows that
\[
\lambda\equiv\frac14(qq+q)\equiv\frac14(q+1)^2-\frac14(q+1)\pmod4
\]
or, since $\frac14(q+1)^2\equiv0\pmod4$, $\lambda\equiv-\frac14(q+1)\equiv\frac14(a-1)\pmod4$. This is precisely the theorem of Article~63 for the case $b=0$.

\SourceLeaf{150}%
